\chapter*{Abstract}
\addcontentsline{toc}{chapter}{Abstract}

System log streams represent the primary telemetry modality for tracking operational health, service reliability, and security integrity in large-scale distributed computing environments and mission-critical networked infrastructures. Modern telemetry streams generate millions of semi-structured log events across massively parallel nodes, making manual diagnostic inspection impossible and demanding automated, real-time anomaly detection pipelines. 

In this capstone research project, we present an exhaustive empirical investigation into the spectrum of log anomaly detection architectures, spanning lightweight linear subspace projections, robust decomposition techniques, and deep semantic reconstruction models. We begin by reproducing the foundational linear baselines established in recent literature (F{\"a}lt et al., IEEE ICSRS 2025) on the structured Hadoop Distributed File System (\textbf{HDFS\_v1}) benchmark (575,061 total blocks, 16,838 anomalous). Under the standard 5-fold cross-validation protocol (25,000 normal and 1,000 anomalous traces; 5,000 normals for fitting), we uncover that an undocumented $L_2$-normalization of raw event count vectors is strictly required to reproduce the published Principal Component Analysis (PCA, $k=6$) ceiling: our implementation achieves $94.60\%$ $F_1$-score, matching the reported $94.64\%$. Furthermore, replacing the literature's test-set oracle percentile cutoff with label-free Extreme Value Theory Peaks-Over-Threshold (POT) thresholding yields $95.95\% \pm 1.47\%$ $F_1$, successfully eliminating test-set leakage while remaining within statistical baseline noise. Conversely, extensive evaluation over 30 random seeds demonstrates that the published Robust PCA (RPCA) claim of $90.55\%$ $F_1$ is irreproducible under strict evaluation protocols (peaking at $87.38\%$), as clean training partitions cause low-rank matrix collapse to rank-2.

Crucially, when extending linear count-based models to the highly interleaved Blue Gene/L (\textbf{BGL}) supercomputing dataset (4,747,963 log lines, 1,848 Drain templates), linear models hit a strict mathematical ceiling of $\sim 65\%$ $F_1$. To overcome this barrier, we engineer a spatial-temporal sequence architecture that groups logs by physical node identity over 6-hour windows (yielding 1,286,817 sequences), projects events into continuous vector spaces via unsupervised Skip-Gram Word2Vec (32-d), aggregates representations using concatenated Mean + Max dual pooling ($z \in \mathbb{R}^{64}$), and reconstructs features through a deep Multi-Layer Perceptron (MLP) Autoencoder ($64 \rightarrow 32 \rightarrow 16 \rightarrow 32 \rightarrow 64$).

To establish unquestionable scientific rigor, we deploy a leak-free 10-fold cross-validation protocol under natural operating prevalence ($10.07\%$ anomalies, where every anomaly is tested exactly once and models are trained strictly on healthy behavior with zero test-label feedback). Under pure unsupervised EVT/POT thresholding on validation normal reconstruction errors (\textbf{Mode A}, headline result, $q=0.90$, target $\text{FPR}=10^{-3}$, $0/10$ fallbacks triggered), our system attains:
\begin{itemize}
    \item \textbf{Precision:} $93.76\% \pm 2.54\%$
    \item \textbf{Recall:} $82.61\% \pm 1.96\%$
    \item \textbf{F1-Score:} $\mathbf{87.80\% \pm 1.31\%}$
\end{itemize}
Under semi-supervised validation tuning (\textbf{Mode B}), the system attains Precision $93.06\% \pm 2.59\%$, Recall $84.71\% \pm 1.19\%$, and $F_1$-score $\mathbf{88.66\% \pm 1.07\%}$. Threshold-free ranking metrics confirm exceptional discriminative power, achieving Precision-Recall AUC of $\mathbf{91.05\% \pm 0.32\%}$ and ROC-AUC of $\mathbf{97.34\% \pm 0.15\%}$.

We definitively deconstruct earlier literature claims of $\sim 90.44\%$ $F_1$ by evaluating our exact trained models on an artificially balanced test set ($50\%$ anomalies): our models score $\mathbf{90.16\%}$ (Mode A) and $\mathbf{91.36\%}$ (Mode B) $F_1$. We prove that the gap between $87.80\%$ and $90.44\%$ is fundamentally driven by base-rate prevalence and test contamination in legacy scripts rather than algorithmic regression. Architectural ablations demonstrate that the non-linear Autoencoder significantly outperforms linear PCA ($59.97\%$ Mode A $F_1$) and Isolation Forest ($2.07\%$ Mode A $F_1$).

Furthermore, temporal robustness tests across a 70/10/20 chronological split expose a critical production challenge: Mode A $F_1$ collapses to $0.03\%$ due to severe concept drift and an Out-of-Vocabulary (OOV) token rate of $68.91\%$, proving that static embeddings cannot survive system evolution without periodic retraining. Finally, hardware benchmarks on an Apple M2 CPU reveal that the serialized autoencoder requires only \textbf{5,328 parameters} ($70.25\text{ KB}$ footprint), achieves a single-window latency of \textbf{43.2 microseconds} (batch latency $0.37\text{ }\mu\text{s/window}$), and delivers an end-to-end throughput of \textbf{62,076 windows/second} without GPU acceleration. We synthesize these empirical discoveries into a universal \textit{Dataset-Dependency Rule} governing telemetry anomaly detection systems.

\vspace{0.8cm}
\noindent\textbf{Keywords:} Log Anomaly Detection, System Telemetry, Autoencoder, Word2Vec, Extreme Value Theory, Peaks-Over-Threshold, Base-Rate Prevalence, Model Efficiency.

